% GENERATED FILE. Edit canonical lessons, then run npm run generate:books.
% canonical-source-hash: fnv1a64:e541ddb93eb50818
% canonical-lessons: ML-C350-talakkettu, ML-C350-patradhipar, ML-C350-patippu, ML-C350-vayanakkaran, ML-C350-accatikkuka

\chapter{Headlines and Editions}
\label{ch:ml-headlines-and-editions}
\hlchaptermodality{350}

\begin{chapteropening}
\textbf{By the end of this chapter:} \emph{I can say a headline, an editor, an edition, a reader and to print.}

Complete the last lesson of chapter 350: I can say a headline, an editor, an edition, a reader and to print.
\end{chapteropening}

\section[\texorpdfstring{\ml{തലക്കെട്ട്} (talakkeṭṭŭ)}{തലക്കെട്ട് (talakkeṭṭŭ)}]{\ml{തലക്കെട്ട്} (talakkeṭṭŭ) — a headline; a title}

\label{lesson:ML-C350-talakkettu}



\begin{quote}
Before the new one: say the Malayalam for to happen, then the Malayalam for to announce.
\end{quote}

\subsection*{You'll want to know: \ml{തലക്കെട്ട്}}
\textbf{\ml{തലക്കെട്ട്}} — \emph{talakkeṭṭŭ} — \textquotedblleft{}a headline; a title\textquotedblright{}.

\begin{quote}
\ml{പത്രത്തിലെ} \ml{വലിയ} \ml{തലക്കെട്ട്} \ml{ഞാൻ} \ml{വായിച്ചു}.
\end{quote}

\begin{quote}
\emph{patrattile valiya talakkeṭṭŭ ñān vāyiccu.}
\end{quote}

\begin{quote}
I read the big headline in the newspaper.
\end{quote}

Literally a head-band: the line tied across the top of a story or a chapter.

\begin{cousinweb}
\textbf{\ml{തലക്കെട്ട്}} (\emph{talakkeṭṭŭ}) is a native Dravidian compound of \ml{തല} (\emph{tala}), \textquotedblleft{}head\textquotedblright{} (DEDR 3103), and \ml{കെട്ട്} (\emph{keṭṭŭ}), \textquotedblleft{}a tying, a band\textquotedblright{}, the noun of \ml{കെട്ടുക}, \textquotedblleft{}to tie\textquotedblright{} (DEDR 1147).
\end{cousinweb}

\begin{grammarlens}[title={What you've built}]
One more word for telling what happened.
\end{grammarlens}

\subsection*{Your turn}
\begin{itemize}
\raggedright
  \item \emph{Say it:} \emph{talakkeṭṭŭ}
  \item \emph{Say it:} \emph{talakkeṭṭŭ}, once more
  \item \emph{Say it:} \emph{patrattile valiya talakkeṭṭŭ ñān vāyiccu.} — I read the big headline in the newspaper.
  \item \emph{Recall:} say the Malayalam for to happen, then the Malayalam for to announce, then say \emph{talakkeṭṭŭ} again
\end{itemize}

\begin{tcolorbox}[breakable,colback=teal!4,colframe=teal!35!black,title={Before you move on}]
What does \ml{തലക്കെട്ട്} mean? (\textquotedblleft{}A headline; a title\textquotedblright{}.) Say it once more.
\end{tcolorbox}
\section[\texorpdfstring{\ml{പത്രാധിപർ} (patrādhipar)}{പത്രാധിപർ (patrādhipar)}]{\ml{പത്രാധിപർ} (patrādhipar) — an editor}

\label{lesson:ML-C350-patradhipar}



\begin{quote}
Before the new one: say the Malayalam for to announce, then the Malayalam for a headline.
\end{quote}

\subsection*{You'll want to know: \ml{പത്രാധിപർ}}
\textbf{\ml{പത്രാധിപർ}} — \emph{patrādhipar} — \textquotedblleft{}an editor\textquotedblright{}.

\begin{quote}
\ml{പത്രാധിപർ} \ml{ലേഖകന്റെ} \ml{റിപ്പോർട്ട്} \ml{വായിച്ചു}.
\end{quote}

\begin{quote}
\emph{patrādhipar lēkhakanṟe ṟippōrṭṭŭ vāyiccu.}
\end{quote}

\begin{quote}
The editor read the correspondent's report.
\end{quote}

The -\ml{ർ} ending marks respect, so one editor is \ml{പത്രാധിപർ}; the plainer \ml{പത്രാധിപൻ} is heard too.

\begin{cousinweb}
\textbf{\ml{പത്രാധിപർ}} (\emph{patrādhipar}) is built from two Sanskrit words: \emph{patra}, \textquotedblleft{}a leaf, a written sheet\textquotedblright{}, which gave \ml{പത്രം} (\emph{patraṁ}), \textquotedblleft{}a newspaper\textquotedblright{}, and \emph{adhipa}, \textquotedblleft{}a ruler, a commander\textquotedblright{} (Monier-Williams). The editor is the chief of the paper.
\end{cousinweb}

\begin{grammarlens}[title={What you've built}]
One more word for telling what happened.
\end{grammarlens}

\subsection*{Your turn}
\begin{itemize}
\raggedright
  \item \emph{Say it:} \emph{patrādhipar}
  \item \emph{Say it:} \emph{patrādhipar}, once more
  \item \emph{Say it:} \emph{patrādhipar lēkhakanṟe ṟippōrṭṭŭ vāyiccu.} — The editor read the correspondent's report.
  \item \emph{Recall:} say the Malayalam for to announce, then the Malayalam for a headline, then say \emph{patrādhipar} again
\end{itemize}

\begin{tcolorbox}[breakable,colback=teal!4,colframe=teal!35!black,title={Before you move on}]
What does \ml{പത്രാധിപർ} mean? (\textquotedblleft{}An editor\textquotedblright{}.) Say it once more.
\end{tcolorbox}
\section[\texorpdfstring{\ml{പതിപ്പ്} (patippŭ)}{പതിപ്പ് (patippŭ)}]{\ml{പതിപ്പ്} (patippŭ) — an edition; an issue}

\label{lesson:ML-C350-patippu}



\begin{quote}
Before the new one: say the Malayalam for a headline, then the Malayalam for an editor.
\end{quote}

\subsection*{You'll want to know: \ml{പതിപ്പ്}}
\textbf{\ml{പതിപ്പ്}} — \emph{patippŭ} — \textquotedblleft{}an edition; an issue\textquotedblright{}.

\begin{quote}
\ml{ആ} \ml{പുസ്തകത്തിന്റെ} \ml{പുതിയ} \ml{പതിപ്പ്} \ml{ഇന്നലെ} \ml{വന്നു}.
\end{quote}

\begin{quote}
\emph{ā pustakattinṟe putiya patippŭ innale vannu.}
\end{quote}

\begin{quote}
A new edition of that book came out yesterday.
\end{quote}

Also one issue of a paper or a magazine: \ml{ഇന്നത്തെ} \ml{പതിപ്പ്}, today's issue.

\begin{cousinweb}
\textbf{\ml{പതിപ്പ്}} (\emph{patippŭ}) is native Dravidian, a noun built on \ml{പതിക്കുക} (\emph{patikkuka}), \textquotedblleft{}to imprint, to stamp\textquotedblright{} (DEDR 3911), with the noun ending -\ml{പ്പ്}. An edition is one pressing of the type.
\end{cousinweb}

\begin{grammarlens}[title={What you've built}]
One more word for telling what happened.
\end{grammarlens}

\subsection*{Your turn}
\begin{itemize}
\raggedright
  \item \emph{Say it:} \emph{patippŭ}
  \item \emph{Say it:} \emph{patippŭ}, once more
  \item \emph{Say it:} \emph{ā pustakattinṟe putiya patippŭ innale vannu.} — A new edition of that book came out yesterday.
  \item \emph{Recall:} say the Malayalam for a headline, then the Malayalam for an editor, then say \emph{patippŭ} again
\end{itemize}

\begin{tcolorbox}[breakable,colback=teal!4,colframe=teal!35!black,title={Before you move on}]
What does \ml{പതിപ്പ്} mean? (\textquotedblleft{}An edition; an issue\textquotedblright{}.) Say it once more.
\end{tcolorbox}
\section[\texorpdfstring{\ml{വായനക്കാരൻ} (vāyanakkāran)}{വായനക്കാരൻ (vāyanakkāran)}]{\ml{വായനക്കാരൻ} (vāyanakkāran) — a reader}

\label{lesson:ML-C350-vayanakkaran}



\begin{quote}
Before the new one: say the Malayalam for an editor, then the Malayalam for an edition.
\end{quote}

\subsection*{You'll want to know: \ml{വായനക്കാരൻ}}
\textbf{\ml{വായനക്കാരൻ}} — \emph{vāyanakkāran} — \textquotedblleft{}a reader\textquotedblright{}.

\begin{quote}
\ml{ആ} \ml{മാസികയ്ക്ക്} \ml{ഒരുപാട്} \ml{വായനക്കാർ} \ml{ഉണ്ട്}.
\end{quote}

\begin{quote}
\emph{ā māsikaykkŭ orupāṭŭ vāyanakkār uṇṭŭ.}
\end{quote}

\begin{quote}
That magazine has a lot of readers.
\end{quote}

Plural: \ml{വായനക്കാർ}, \emph{vāyanakkār}.

A woman reader is \ml{വായനക്കാരി}. The -\ml{ക്കാരൻ} ending makes a person from an activity, as in \ml{കച്ചവടക്കാരൻ}.

\begin{cousinweb}
\textbf{\ml{വായനക്കാരൻ}} (\emph{vāyanakkāran}) is a Malayalam compound of \ml{വായന} (\emph{vāyana}), \textquotedblleft{}reading\textquotedblright{}, the noun of \ml{വായിക്കുക}, \textquotedblleft{}to read\textquotedblright{}, and the ending -\ml{ക്കാരൻ}, \textquotedblleft{}one who does it\textquotedblright{}. The verb is usually traced to Sanskrit \emph{vāc}, \textquotedblleft{}speech\textquotedblright{}: reading was first reading aloud.
\end{cousinweb}

\begin{grammarlens}[title={What you've built}]
One more word for telling what happened.
\end{grammarlens}

\subsection*{Your turn}
\begin{itemize}
\raggedright
  \item \emph{Say it:} \emph{vāyanakkāran}
  \item \emph{Say it:} \emph{vāyanakkāran}, once more
  \item \emph{Say it:} \emph{ā māsikaykkŭ orupāṭŭ vāyanakkār uṇṭŭ.} — That magazine has a lot of readers.
  \item \emph{Recall:} say the Malayalam for an editor, then the Malayalam for an edition, then say \emph{vāyanakkāran} again
\end{itemize}

\begin{tcolorbox}[breakable,colback=teal!4,colframe=teal!35!black,title={Before you move on}]
What does \ml{വായനക്കാരൻ} mean? (\textquotedblleft{}A reader\textquotedblright{}.) Say it once more.
\end{tcolorbox}
\section[\texorpdfstring{\ml{അച്ചടിക്കുക} (accaṭikkuka)}{അച്ചടിക്കുക (accaṭikkuka)}]{\ml{അച്ചടിക്കുക} (accaṭikkuka) — to print}

\label{lesson:ML-C350-accatikkuka}



\begin{quote}
Before the new one: say the Malayalam for an edition, then the Malayalam for a reader.
\end{quote}

\subsection*{You'll want to know: \ml{അച്ചടിക്കുക}}
\textbf{\ml{അച്ചടിക്കുക}} — \emph{accaṭikkuka} — \textquotedblleft{}to print\textquotedblright{}.

\begin{quote}
\ml{അവർ} \ml{രാത്രി} \ml{മുഴുവൻ} \ml{പത്രം} \ml{അച്ചടിച്ചു}.
\end{quote}

\begin{quote}
\emph{avar rātri muḻuvan patraṁ accaṭiccu.}
\end{quote}

\begin{quote}
They printed the newspaper all night.
\end{quote}

Past: \ml{അച്ചടിച്ചു}, \emph{accaṭiccu}.

The noun is \ml{അച്ചടി}, printing.

\begin{cousinweb}
\textbf{\ml{അച്ചടിക്കുക}} (\emph{accaṭikkuka}) is a native Dravidian compound of \ml{അച്ച്} (\emph{accŭ}), \textquotedblleft{}a mould, a type\textquotedblright{} (DEDR 47), and \ml{അടിക്കുക} (\emph{aṭikkuka}), \textquotedblleft{}to strike, to stamp\textquotedblright{} (DEDR 77): to print is to strike paper with type.
\end{cousinweb}

\begin{grammarlens}[title={What you've built}]
One more word for telling what happened.
\end{grammarlens}

\subsection*{Your turn}
\begin{itemize}
\raggedright
  \item \emph{Say it:} \emph{accaṭikkuka}
  \item \emph{Say it:} \emph{accaṭikkuka}, once more
  \item \emph{Say it:} \emph{avar rātri muḻuvan patraṁ accaṭiccu.} — They printed the newspaper all night.
  \item \emph{Recall:} say the Malayalam for an edition, then the Malayalam for a reader, then say \emph{accaṭikkuka} again
\end{itemize}

\begin{tcolorbox}[breakable,colback=teal!4,colframe=teal!35!black,title={Before you move on}]
What does \ml{അച്ചടിക്കുക} mean? (\textquotedblleft{}To print\textquotedblright{}.) Say it once more.
\end{tcolorbox}
